
\subsubsection{2.1 ML Reproducibility}

Raff [2019] provided the foundational labeled corpus for ML reproducibility research: 255 top-venue papers (NeurIPS, ICML, ICLR, JMLR) from 1984--2017, each with a binary reproducibility label from manual re-implementation attempts by a single researcher. The corpus established that approximately 50.8\% of papers succeed under independent reproduction and identified paper-level features predictive of success: pseudocode presence, equation count, and hyperparameter reporting. Critically, Raff analyzed \textit{paper-quality} features rather than \textit{dataset-quality} features --- the dataset-level signals this study investigates were not available from public APIs at the time of his study, and Raff's original feature set does not include metadata-observable misuse signals.

The ML Reproducibility Challenge \citep{Pineau2021Improving} extended this work to post-2019 papers using multi-annotator reproduction attempts. NeurIPS, ICML, and other venues have adopted reproducibility checklists requiring code submission, dataset specification, and hyperparameter reporting. These interventions are normative --- they specify what should be done --- but do not empirically test whether documentation presence predicts reproduction success.

Gundersen and Kjensmo [2018] analyzed ML papers for reproducibility factors including experimental setup documentation. Kapoor and Narayanan [2023] identified data leakage as a systematic driver of reproducibility failure in ML across domains. Neither work connects dataset-level metadata APIs to reproducibility labels.

\subsubsection{2.2 Dataset Documentation}

Gebru et al. [2021] introduced Datasheets for Datasets, proposing a structured documentation framework for ML datasets covering motivation, composition, collection process, uses, and maintenance. The Datasheets schema has been adopted as the basis for HuggingFace Hub dataset cards and is the foundation of the field-presence completeness metric used in this study. Gebru et al. propose the schema normatively: they do not test whether dataset card completeness predicts downstream model failure.

Bender et al. [2021] and Gebru et al. [2021] both argue that documentation scope fields (\texttt{intended\textbackslash \_use}, \texttt{out\textbackslash \_of\textbackslash \_scope\textbackslash \_use}) are critical for preventing out-of-context dataset application --- the mechanism theorized to predict reproducibility failure. The present work is the first attempt to operationalize this theoretical link against a labeled reproducibility outcome dataset, though the operationalization was blocked by the infrastructure gap reported here.

YoungXinyu1802 et al. [2024] analyzed 7,433 HuggingFace Hub dataset cards, finding improving completeness over time and identifying field-level completion rates. Their analysis covers the overall HF Hub ecosystem but does not characterize coverage for specific reproducibility corpora, does not analyze pre-2018 ML benchmark datasets, and does not connect card completeness to reproducibility outcomes. The finding in this study that even found HF cards average only 0.41 mean field-presence score complements their temporal analysis by revealing the thinness of retroactively-created cards for foundational benchmarks.

\subsubsection{2.3 Benchmark Concentration and Underspecification}

D'Amour et al. [2022] introduced underspecification as a challenge for credibility in ML: models trained on the same data to the same accuracy may behave very differently at deployment because the training distribution fails to select among many valid predictors. Highly concentrated benchmark datasets --- where the community has trained and evaluated on the same splits over many years --- accumulate dataset-specific artifacts that models learn. D'Amour et al. do not operationalize concentration as a measurable quantity or link it to Raff's reproducibility labels.

Recht et al. [2019] and Engstrom et al. [2020] demonstrated that models trained on CIFAR-10 and ImageNet show substantial accuracy drops on new test sets, consistent with artifact overfitting. These findings provide empirical grounding for the concentration hypothesis: high usage concentration in a benchmark may signal elevated artifact accumulation risk.

\subsubsection{2.4 Dataset Infrastructure for Reproducibility Studies}

Vanschoren et al. [2014] introduced OpenML as a platform for systematic ML experimentation, providing per-dataset run counts and task creation timestamps. OpenML is designed for tabular, classical ML experiments and has been used in AutoML benchmarking \citep{Feurer2022Auto}. This study is the first to systematically quantify OpenML's coverage for a deep-learning-heavy reproducibility corpus, finding 22\% coverage concentrated in classical tabular and a small number of NLP and image datasets --- consistent with OpenML's documented design scope but previously unquantified for Raff's corpus.

Papers With Code \citep{Stojnic2020Papers} explicitly links ML papers to datasets, code implementations, and benchmark results across all domains. In contrast to HF Hub and OpenML, Papers With Code covers DL vision, NLP, speech, and tabular datasets and is used in reproducibility research.

\subsubsection{2.5 Positioning}

No prior work has: (1) empirically tested whether HF Hub and OpenML provide adequate coverage of the Raff 2019 corpus for metadata-based reproducibility auditing, (2) characterized the domain-temporal selection biases of these APIs for pre-2018 ML benchmarks, or (3) proposed Papers With Code as the appropriate primary API for this class of study. This work addresses all three.

